\documentclass[11pt,a4paper]{article}
\usepackage[utf8]{inputenc}
\usepackage[T1]{fontenc}
\usepackage{geometry}
\usepackage{setspace}
\usepackage{parskip}
\usepackage{booktabs}
\usepackage{longtable}
\usepackage{array}
\usepackage{xcolor}
\usepackage{hyperref}
\usepackage{fancyhdr}
\usepackage{titlesec}

\geometry{left=2cm,right=2cm,top=2.5cm,bottom=2.5cm,headheight=14.5pt}
\onehalfspacing
\hypersetup{colorlinks=true,linkcolor=black,urlcolor=blue,pdftitle={Distribución de Casos de Uso - Sistema Integral de Gestión de Eventos},pdfauthor={Hexacore}}

\definecolor{azul}{RGB}{31,78,121}
\definecolor{gris}{RGB}{242,242,242}
\definecolor{grisoscuro}{RGB}{80,80,80}

\titleformat{\section}{\Large\bfseries\color{azul}}{\thesection.}{0.5em}{}

\pagestyle{fancy}
\fancyhf{}
\fancyhead[L]{Sistema Integral de Gestión de Eventos}
\fancyhead[R]{Distribución de Casos de Uso}
\fancyfoot[C]{\thepage}

% Column types for the use-case table
\newcolumntype{P}[1]{>{\raggedright\arraybackslash}p{#1}}

\begin{document}

\begin{titlepage}
\centering
\vspace*{2.5cm}
{\Huge\bfseries Sistema Integral de Gestión de Eventos\par}
\vspace{0.7cm}
{\Large Arquitectura de Software\par}
\vspace{1.2cm}
{\Large\bfseries Distribución de Casos de Uso por Integrante\par}
\vspace{2cm}
{\large Equipo Hexacore\par}
\vfill
\begin{tabular}{rl}
\textbf{Materia:} & Arquitectura de Software\\
\textbf{Equipo:} & Hexacore\\
\textbf{Integrantes:} & Daniel Cristancho, Samuel Emperador, Sebastián Sánchez, Diego Coronado\\
\textbf{Total de casos de uso:} & 32 (8 por integrante)\\
\textbf{Fecha:} & Agosto de 2026
\end{tabular}
\end{titlepage}

\tableofcontents
\newpage

\section*{Resumen de la distribución}
\addcontentsline{toc}{section}{Resumen de la distribución}

Los 32 casos de uso del proyecto se reparten en bloques de 8 por integrante,
respetando el mínimo de 5 casos de uso por persona exigido por el curso y
manteniendo agrupados los bloques temáticos ya definidos por el equipo
(Boletería, Personal, Pedidos/Alimentos, Logística/Operación, Parqueaderos y
los módulos CRUD administrativos añadidos posteriormente).

\begin{longtable}{P{4.5cm}P{4cm}P{6.5cm}}
\toprule
\textbf{Integrante} & \textbf{Casos de uso} & \textbf{Bloque temático}\\
\midrule
\endhead
Daniel Cristancho & CU-001 -- CU-005, CU-021 -- CU-023 & Boletería / Compra de entradas + Parqueadero (reserva, ingreso/salida, asignación de espacios)\\[0.3em]
Samuel Emperador & CU-006 -- CU-010, CU-024 -- CU-026 & Personal para eventos + Parqueadero (cobro, control de ocupación) + Gestión de Eventos\\[0.3em]
Sebastián Sánchez & CU-011 -- CU-015, CU-027 -- CU-029 & Pedidos y establecimientos de alimentos + Cuentas de usuario, Roles y Permisos, Recintos y Zonas\\[0.3em]
Diego Coronado & CU-016 -- CU-020, CU-030 -- CU-032 & Logística y operación del evento + Proveedores, Pagos/Conciliación, Reportes y Analítica\\
\bottomrule
\caption{Resumen de la distribución de los 32 casos de uso entre los 4 integrantes.}
\end{longtable}

\newpage

%===================== DANIEL CRISTANCHO =====================
\section{Daniel Cristancho}

\begin{longtable}{P{1.6cm}P{4cm}P{9.4cm}}
\toprule
\textbf{Id} & \textbf{Caso de uso} & \textbf{Descripción}\\
\midrule
\endfirsthead
\toprule
\textbf{Id} & \textbf{Caso de uso} & \textbf{Descripción}\\
\midrule
\endhead

CU-001 & Compra de entradas &
El usuario consulta los eventos disponibles, selecciona uno, elige la localidad y cantidad de
entradas, realiza el pago y recibe sus boletos digitales con código QR.\\[0.5em]

CU-002 & Validar QR &
Al llegar al evento, el asistente presenta su código QR. El sistema valida que la entrada sea
auténtica, que no haya sido utilizada, y registra el ingreso.\\[0.5em]

CU-003 & Cancelaciones y devoluciones &
El usuario solicita la cancelación de una compra, se verifican las condiciones del evento y se
procesa el reembolso.\\[0.5em]

CU-004 & Promociones &
Durante el proceso de compra, el cliente aplica un código promocional para obtener un
descuento.\\[0.5em]

CU-005 & Consultar evento &
El usuario puede explorar los eventos filtrando por categoría, fecha, ciudad o artista, y
consultar detalles como horarios, localidades, precios y disponibilidad.\\[0.5em]

CU-021 & Reservar parqueadero &
Permite que el cliente reserve un espacio de parqueadero para el evento antes de su
llegada.\\[0.5em]

CU-022 & Ingreso y salida de vehículos &
Controla el ingreso y la salida de los vehículos al parqueadero del evento.\\[0.5em]

CU-023 & Asignación de espacios &
Asigna automáticamente un espacio disponible dentro de la zona de parqueadero
correspondiente.\\
\bottomrule
\caption{Casos de uso a cargo de Daniel Cristancho.}
\end{longtable}

\newpage

%===================== SAMUEL EMPERADOR =====================
\section{Samuel Emperador}

\begin{longtable}{P{1.6cm}P{4cm}P{9.4cm}}
\toprule
\textbf{Id} & \textbf{Caso de uso} & \textbf{Descripción}\\
\midrule
\endfirsthead
\toprule
\textbf{Id} & \textbf{Caso de uso} & \textbf{Descripción}\\
\midrule
\endhead

CU-006 & Gestión del Mercado Secundario de Entradas &
Permite que un usuario que ya adquirió una entrada la publique para reventa y que otro usuario
la compre de forma segura, garantizando el cambio de propiedad, la invalidación del QR anterior
y la generación de un nuevo QR.\\[0.5em]

CU-007 & Gestión de Personal para Eventos (Administración y Asignación) &
Permite administrar la información del personal que participa en los eventos (crear, consultar,
actualizar y desactivar) y asignarlo a un evento validando su rol, disponibilidad y horario,
evitando conflictos de asignación.\\[0.5em]

CU-008 & Gestionar Cambios de Turno &
Permite que un empleado solicite un cambio de turno y que el supervisor lo apruebe asignando un
reemplazo adecuado, sin generar conflictos de horario.\\[0.5em]

CU-009 & Registrar Ingreso y Salida del Personal &
Controla la asistencia del personal durante el evento mediante el registro de entrada y salida,
y calcula automáticamente las horas trabajadas.\\[0.5em]

CU-010 & Gestionar Evacuación ante Emergencias &
Permite que el personal autorizado active, coordine y finalice un protocolo de evacuación
durante una emergencia, utilizando información en tiempo real sobre personas presentes, zonas
ocupadas, salidas disponibles, rutas de evacuación y personal operativo.\\[0.5em]

CU-024 & Cobro &
Procesa el cobro del servicio de parqueadero según el tiempo de permanencia o la tarifa de
reserva.\\[0.5em]

CU-025 & Control de ocupación &
Supervisa en tiempo real la ocupación de las zonas de parqueadero del evento.\\[0.5em]

CU-026 & Gestión de Eventos (CRUD) &
El organizador crea, edita, publica y cancela eventos, definiendo su información general, tipos
de entrada, precios por zona y aforo máximo.\\
\bottomrule
\caption{Casos de uso a cargo de Samuel Emperador.}
\end{longtable}

\newpage

%===================== SEBASTIAN SANCHEZ =====================
\section{Sebastián Sánchez}

\begin{longtable}{P{1.6cm}P{4cm}P{9.4cm}}
\toprule
\textbf{Id} & \textbf{Caso de uso} & \textbf{Descripción}\\
\midrule
\endfirsthead
\toprule
\textbf{Id} & \textbf{Caso de uso} & \textbf{Descripción}\\
\midrule
\endhead

CU-011 & Gestión integral de pedidos &
Permite al cliente realizar un pedido de alimentos durante un evento, desde la selección del
establecimiento y los productos hasta el pago, registro del pedido y generación del código
QR.\\[0.5em]

CU-012 & Administrar menú e inventario &
Permite al proveedor administrar los productos de su menú y su disponibilidad para la venta
durante los eventos.\\[0.5em]

CU-013 & Gestionar preparación y entrega del pedido &
Permite al establecimiento gestionar un pedido confirmado desde su recepción y preparación hasta
la validación del cliente y la entrega final, garantizando que el pedido correcto sea entregado
una sola vez.\\[0.5em]

CU-014 & Validar y entregar pedido &
Valida el código QR presentado por el cliente y autoriza la entrega únicamente cuando el pedido
exista, esté pagado, esté listo y no haya sido entregado previamente.\\[0.5em]

CU-015 & Gestionar establecimientos de alimentos &
Permite al organizador o administrador registrar, consultar, actualizar, habilitar o deshabilitar
los establecimientos de alimentos asociados a un evento, definiendo su información operativa y
su punto de entrega.\\[0.5em]

CU-027 & Gestión de Cuentas de Usuario &
Cualquier usuario del sistema se registra, inicia sesión, recupera su contraseña y edita su
perfil; el administrador puede además activar, desactivar o eliminar cuentas.\\[0.5em]

CU-028 & Gestión de Roles y Permisos &
El administrador crea, edita y elimina roles del sistema, y asigna o revoca permisos sobre cada
módulo (eventos, personal, parqueaderos, alimentos, reportes).\\[0.5em]

CU-029 & Gestión de Recintos y Zonas &
Registra y mantiene los recintos disponibles, sus zonas, capacidad por zona y mapa de
distribución, como insumo para la creación de eventos y la venta de entradas por zona.\\
\bottomrule
\caption{Casos de uso a cargo de Sebastián Sánchez.}
\end{longtable}

\newpage

%===================== DIEGO CORONADO =====================
\section{Diego Coronado}

\begin{longtable}{P{1.6cm}P{4cm}P{9.4cm}}
\toprule
\textbf{Id} & \textbf{Caso de uso} & \textbf{Descripción}\\
\midrule
\endfirsthead
\toprule
\textbf{Id} & \textbf{Caso de uso} & \textbf{Descripción}\\
\midrule
\endhead

CU-016 & Planificar logística del evento &
El coordinador logístico organiza todos los recursos, personal y proveedores necesarios para
garantizar el correcto desarrollo del evento.\\[0.5em]

CU-017 & Asignar personal operativo &
El coordinador logístico distribuye el personal disponible en las diferentes zonas del
evento.\\[0.5em]

CU-018 & Gestionar turno y asistencia del personal &
Permite que el personal operativo solicite cambios de turno y registre su ingreso y salida real
durante el evento.\\[0.5em]

CU-019 & Monitorear el evento &
Supervisa el desarrollo del evento en tiempo real.\\[0.5em]

CU-020 & Gestionar incidentes &
Registra y atiende cualquier incidente ocurrido durante el evento.\\[0.5em]

CU-030 & Gestión de Proveedores &
Registra, edita y da de baja proveedores externos (catering, sonido, seguridad privada, etc.)
que pueden ser asignados en la planificación logística de un evento.\\[0.5em]

CU-031 & Gestión de Pagos y Conciliación Financiera &
El administrador financiero supervisa, concilia y audita todos los pagos procesados por el
sistema (entradas, mercado secundario, parqueaderos, alimentos), y gestiona reembolsos
administrativos cuando el proceso automático de cancelación no aplica.\\[0.5em]

CU-032 & Gestión de Reportes y Analítica del Evento &
El organizador y el administrador consultan reportes y dashboards con métricas del evento
(ventas, ocupación, ingresos, personal, incidentes) para tomar decisiones antes, durante y
después del evento.\\
\bottomrule
\caption{Casos de uso a cargo de Diego Coronado.}
\end{longtable}

\end{document}
